\section{Reduction to rectifiable minimizers}
\label{sec:initial-reduction}

The admissible class in the problem does not assume that the closed set is
rectifiable or that the function is bounded.  We establish both properties
locally before using the classical regularity theory.  The essential point is
that a portion of the closed set which carries no jump can be removed by a
Dirichlet replacement, with an arbitrarily small error on the boundary of the
replacement disk.

We use the following standard notation.  A function is in $BV(W)$ if its
distributional derivative is a finite vector-valued Radon measure on $W$.
Its derivative decomposes into an absolutely continuous part, a jump part,
and a Cantor part.  The jump part is concentrated on the approximate jump
set $J_u$, where two distinct one-sided approximate limits exist across a
line.  The space $SBV(W)$ consists of the $BV$ functions whose Cantor part
vanishes.  We use two basic facts from the $BV$ structure theorem: $J_u$ is
countably $1$-rectifiable (covered, up to a set of zero length, by
countably many Lipschitz images of intervals), and the Cantor part gives zero mass to every set
of $\sigma$-finite $\mathcal H^1$ measure
\cite[Theorem~3.78 and Proposition~3.92(c)]{ambrosio_fusco_pallara2000}. The same notation with the
subscript ${\rm loc}$ requires these properties on relatively compact
open subsets.

\begin{proposition}[Reduction to a rectifiable pair]
\label{prop:initial-reduction}
Let $\Omega\subset\mathbb R^2$ be bounded and open, let
$g\in L^\infty(\Omega)$, and let $(u,K)$ be an admissible absolute
minimizer for $\mathcal F_g$ as defined in Section~\ref{sec:problem}.  In particular,
$K$ is relatively closed, $\mathcal H^1(K)<\infty$, and
$u\in L^2(\Omega)\cap W^{1,2}(\Omega\setminus K)$; no rectifiability
of $K$ is assumed.  Then
\[
 u\in L^\infty_{\rm loc}(\Omega)\cap SBV_{\rm loc}(\Omega),
 \qquad
 \mathcal H^1(K\setminus J_u)=0.
\]
The function satisfies $\Delta u=u-g$ on $\Omega\setminus K$ and has
there a $C^{1,\alpha}_{\rm loc}$ representative for every $\alpha\in(0,1)$.
If the pair is reduced, then
\begin{equation}
 K=\operatorname{spt}_{\Omega}(\mathcal H^1\!\llcorner J_u)
   =\overline{J_u}^{\,\Omega}.
 \label{eq:jump-support}
\end{equation}
Consequently, in the reduced case, on every smooth open set $W\Subset\Omega$, the restricted
pair is a reduced rectifiable absolute minimizer with bounded function.
\end{proposition}

The proof uses a standard Dirichlet theorem in the $SBV$ formulation.
We record the precise version needed here. Its weak minimizer is an
auxiliary replacement: essential closure will turn its jump set into an
admissible closed comparison set, allowing us to test the original pair.
For a disk $D$ and
$w\in C^1_c(\mathbb R^2)$, consider the functional
\begin{equation}
 \mathcal J_D(z)
 :=\int_D\bigl(|\nabla z|^2+|z-g|^2\bigr)\,dx
      +\mathcal H^1(J_z\cap\overline D)
 \label{eq:weak-dirichlet-energy}
\end{equation}
on $z\in SBV(\mathbb R^2)$ satisfying $z=w$ almost everywhere on
$\mathbb R^2\setminus D$.  Here $J_z$ is the jump set of the extension
to the whole plane.  In particular, a mismatch of boundary traces is
charged in \eqref{eq:weak-dirichlet-energy}.

\begin{lemma}[The smooth-data Dirichlet input]
\label{lem:dirichlet-input}
Let $D\subset\mathbb R^2$ be a disk, $g\in L^\infty(D)$, and
$w\in C^1_c(\mathbb R^2)$.  The problem
\eqref{eq:weak-dirichlet-energy} has a bounded minimizer $z$, and
\begin{equation}
 \mathcal H^1\bigl((\overline{J_z}\setminus J_z)
                  \cap\overline D\bigr)=0.
 \label{eq:dirichlet-essential-closure}
\end{equation}
In particular, $C:=\overline{J_z}\cap\overline D$ is compact and has
finite length, and $z$ belongs to $W^{1,2}_{\rm loc}$ on the complement
of $\overline{J_z}$.  On each relatively open portion of $\partial D$ disjoint from $C$,
its interior and exterior Sobolev traces agree almost everywhere.
\end{lemma}

This is the smooth-boundary, smooth-exterior-data case of the Dirichlet
existence and essential-closure theorem; see
\cite[Appendix~B]{delellis_focardi2025}, specifically
Theorem~B.3.1(c), Corollary~B.3.2(b), and Proposition~B.4.2(c).
The final two assertions also follow directly from the $SBV$ decomposition:
on an open set which misses the closed jump set there is no singular
derivative, while the absolutely continuous gradient is square integrable.
Across a jump-free boundary portion, apply this observation to the
whole-plane extension $z$, whose exterior value is $w$.

We explain explicitly why the boundary in
\eqref{eq:dirichlet-essential-closure} introduces no extra length.
The boundary density theorem supplies a positive lower bound for
$\mathcal H^1(J_z\cap B_r(x))/r$ as $r\downarrow0$ at points
$x\in\overline{J_z}\cap\partial D$.
Put $\mu=\mathcal H^1\!\llcorner J_z$ and
$\nu=\mathcal H^1\!\llcorner\partial D$.
The restriction of $\mu$ to the circle is
$\nu\!\llcorner(J_z\cap\partial D)$, while its remaining part is
singular with respect to $\nu$.  Differentiation of measures, together
with $\nu(B_r(x))/(2r)\to1$, gives
\[
 \frac{\mu(B_r(x))}{r}\longrightarrow0
 \quad\text{for $\mathcal H^1$-almost every }
 x\in\partial D\setminus J_z.
\]
This contradicts the boundary lower bound at any such point in the
closed jump set.  Hence the boundary portion of the difference in
\eqref{eq:dirichlet-essential-closure} is null.  The interior portion
is the usual interior essential-closure theorem.  No $C^1$ matching
of normal derivatives across $\partial D$ is used.

\begin{proof}[Proof of Proposition~\ref{prop:initial-reduction}]
We divide the argument into local boundedness, the $SBV$ extension,
and removal of excess length.

\smallskip
\noindent\emph{Step 1: circles with bounded traces.}
Variations $u+t\varphi$, with
$\varphi\in C^\infty_c(\Omega\setminus K)$, preserve the admissible
class and have compact support.  Their first variation gives
\[
 \int_{\Omega\setminus K}\nabla u\cdot\nabla\varphi
       +(u-g)\varphi\,dx=0.
\]
Thus $\Delta u=u-g$ off $K$.  In two dimensions, local
$W^{1,2}$ membership implies local $L^p$ membership for every finite
$p$.  Interior elliptic estimates therefore give
$u\in W^{2,p}_{\rm loc}(\Omega\setminus K)$ for all such $p$,
and hence the asserted $C^{1,\alpha}_{\rm loc}$ regularity.  We use
this representative off $K$ from now on.

Fix $q\in\Omega$.  For almost every radius $s$ with
$\overline B_s(q)\subset\Omega$, the following two properties hold:
\begin{equation}
 \Sigma_s:=K\cap\partial B_s(q)\text{ is finite},
 \qquad
 \int_{\partial B_s(q)\setminus K}
           (|u|^2+|\nabla u|^2)\,d\mathcal H^1<\infty.
 \label{eq:good-circle}
\end{equation}
The first assertion is Eilenberg's coarea inequality applied to the
$1$-Lipschitz distance map on a Borel set of finite $\mathcal H^1$
measure.  It requires no rectifiability.  The second follows from
polar integration.  To recall the elementary content of the first
inequality, cover a compact portion of $K$ by sets of diameter less
than $\delta$.  On each distance level, every $\delta$-separated
collection contains at most one point in each covering set.  The
image of a covering set lies in an interval no longer than its
diameter.  Integration over the levels, followed by refinement of
the covers and $\delta\downarrow0$, bounds the integral of the
level cardinality by the length of the covered set.

Call a disk satisfying \eqref{eq:good-circle} a good disk.
The circle minus $\Sigma_s$ consists of finitely many open arcs.
On each arc the restriction of $u$ is $C^1$, and its tangential
derivative is square integrable by \eqref{eq:good-circle}.
It is therefore a one-dimensional $W^{1,2}$ function and has finite
one-sided limits at the arc endpoints.  Its supremum is finite.
Taking the maximum over the finitely many arcs shows that the
circle trace $f$ is bounded.  If $\Sigma_s$ is empty, the same
conclusion follows from $W^{1,2}$ on the full circle.

\smallskip
\noindent\emph{Step 2: local boundedness by clipping.}
Fix a good disk $B=B_s(q)$ and choose
\[M>\max\{\|f\|_\infty,\|g\|_\infty\}.\]
Let $T_M(a)=\max\{-M,\min\{a,M\}\}$ and set
\[
 \widetilde u=T_M(u)\quad\text{in }B,
 \qquad \widetilde u=u\quad\text{in }\Omega\setminus B.
\]
The traces agree on $\partial B\setminus K$.  Sobolev gluing there,
together with the chain rule for $T_M$, shows that
$\widetilde u\in W^{1,2}(\Omega\setminus K)\cap L^2(\Omega)$.
Its difference from $u$ is supported in $\overline B$.
Choose an open $V$ with
$\overline B\Subset V\Subset\Omega$ and use
$(\widetilde u,K)$ in the defining comparison on $V$.
The Dirichlet term does not increase, and the fidelity term
strictly decreases wherever $|u|>M$ in $B$.
For example, on $\{u>M\}$ the decrease is
$(u-M)(u+M-2g)>0$.
Minimality therefore gives $|u|\le M$ almost everywhere in $B$.
Good disks exist with arbitrarily small radii about every point.
A finite cover of each compact subset of $\Omega$ proves
$u\in L^\infty_{\rm loc}(\Omega)$.

\smallskip
\noindent\emph{Step 3: the $SBV$ extension.}
We give the finite-ball extension argument underlying
\cite[Proposition~4.4]{ambrosio_fusco_pallara2000}; it requires no
rectifiability of the closed set.
Choose smooth open sets $W\Subset W'\Subset\Omega$.
The definition of $\mathcal H^1$ and compactness give finite
covers of $K\cap\overline W$ by open balls of radii $a_{ij}$ such
that
\[
 \max_i a_{ij}\longrightarrow0,
 \qquad \sum_i a_{ij}\le C(1+\mathcal H^1(K)),
 \qquad \overline{G_j}\subset W',
 \quad G_j:=\bigcup_i B_{a_{ij}}.
\]
Their areas tend to zero and their perimeters satisfy
$\operatorname{Per}(G_j)\le2\pi\sum_i a_{ij}$.
Moreover, $\partial G_j\cap W$ misses $K$, since the open union
covers $K\cap\overline W$.
Set $u_j=0$ on $G_j\cap W$ and $u_j=u$ elsewhere in $W$.
Integration by parts along the exposed circular arcs gives
\[
 |Du_j|(W)
 \le\int_{W\setminus K}|\nabla u|\,dx
      +2\pi\|u\|_{L^\infty(W')}\sum_i a_{ij}.
\]
This estimate can equivalently be read as the finite-perimeter
product rule.  Test fields compactly supported in $W$ produce no
term on $\partial W$.
Since $u_j\to u$ in $L^1(W)$, lower semicontinuity proves
$u\in BV(W)$.

Off $K$ the derivative is absolutely continuous, so its Cantor
part is supported on $K$.  That part vanishes on sets of finite
$\mathcal H^1$ measure, and hence it vanishes everywhere in $W$.
Thus $u\in SBV_{\rm loc}(\Omega)$.
Continuity off $K$ gives $J_u\subset K$.
Finally, $K$ has area zero, so the absolutely continuous gradient
in the $SBV$ decomposition agrees almost everywhere with the
original off-set gradient.

We have now placed the function in the weak admissible class.
It remains to show that the original closed set has exactly the
length charged by the jump set.

\smallskip
\noindent\emph{Step 4: smooth boundary data and Dirichlet replacement.}
Fix a good disk $D=B_s(q)$ and its trace $f$.
Choose closed finite unions of circle arcs
$E_m\subset\partial D$ containing the finite set
$\Sigma_s=K\cap\partial D$ in their relative interiors, with
\[
 \ell_m:=\mathcal H^1(E_m)\longrightarrow0.
\]
There are $C^1$ functions $a_m$ on the circle which agree with
$f$ outside $E_m$: on each short arc interpolate the endpoint
values and first tangential derivatives, and keep $f$ on the
remaining arcs.  A radial collar and cutoff extend $a_m$ to a
function $w_m\in C^1_c(\mathbb R^2)$.
Each such extension has bounded first derivatives; no bound
uniform in $m$ is needed.
Let $z_m$ be the minimizer in
Lemma~\ref{lem:dirichlet-input} for exterior datum $w_m$.

The function equal to $u$ in $D$ and to $w_m$ outside $D$ is in
$SBV(\mathbb R^2)$ by the trace gluing formula
\cite[Theorems~3.84 and~3.87, Corollary~3.89]{ambrosio_fusco_pallara2000}.
Its jump set on the circle is contained, up to a null set, in
$E_m$.  Consequently,
\begin{equation}
 \mathcal J_D(z_m)
 \le \int_D\bigl(|\nabla u|^2+|u-g|^2\bigr)\,dx
       +\mathcal H^1(J_u\cap D)+\ell_m.
 \label{eq:dirichlet-replacement-bound}
\end{equation}
Put
\[
 C_m=\overline{J_{z_m}}\cap\overline D,
 \qquad
 L_m=(K\setminus D)\cup C_m\cup E_m,
 \qquad
 v_m=\begin{cases}z_m&\text{in }D,\\u&\text{in }\Omega\setminus D.
       \end{cases}
\]
The set $L_m$ is relatively closed and has finite length by
\eqref{eq:dirichlet-essential-closure}.
The function $v_m$ is locally $W^{1,2}$ away from $L_m$.
Near a circle point outside $L_m$, the interior trace of $z_m$
matches $w_m$ almost everywhere by Lemma~\ref{lem:dirichlet-input},
and that datum matches $u$ because this boundary portion lies
outside $E_m$.
This verifies Sobolev gluing at every possible seam.
The function and its gradient have finite global $L^2$ norms,
so $v_m\in W^{1,2}(\Omega\setminus L_m)\cap L^2(\Omega)$.

Both $L_m\mathbin\Delta K$ and $\operatorname{spt}(v_m-u)$ lie
in $\overline D$.
We may therefore compare in any open
$V$ with $\overline D\Subset V\Subset\Omega$.
Outside $D$ the volume terms agree, and the unchanged portion
of $K$ has the same length.
Since $K\cap\partial D$ is finite, it has zero length.
Using \eqref{eq:dirichlet-essential-closure}, minimality and
\eqref{eq:dirichlet-replacement-bound} give
\begin{align*}
 \mathcal H^1(K\cap D)
 &\le \mathcal J_D(z_m)+\ell_m
       -\int_D\bigl(|\nabla u|^2+|u-g|^2\bigr)\,dx\\
 &\le \mathcal H^1(J_u\cap D)+2\ell_m.
\end{align*}
The two boundary errors have distinct sources: one charges
the mismatched weak trace, and the other inserts $E_m$ in the
closed comparison set.
Letting $m\to\infty$ and using $J_u\subset K$ proves
$\mathcal H^1((K\setminus J_u)\cap D)=0$.
A countable good-disk cover of $\Omega$ yields
$\mathcal H^1(K\setminus J_u)=0$.

\smallskip
\noindent\emph{Step 5: reduced support and restriction.}
If $K$ is reduced, equality of the two length measures gives
\[
 K=\operatorname{spt}_{\Omega}(\mathcal H^1\!\llcorner K)
   =\operatorname{spt}_{\Omega}(\mathcal H^1\!\llcorner J_u).
\]
The support of the jump-length measure is contained in the relative
closure of $J_u$, while $J_u\subset K$ and relative closedness give the
reverse containment in $K$. Thus
\[
 K=\operatorname{spt}_{\Omega}(\mathcal H^1\!\llcorner J_u)
 \subset\overline{J_u}^{\,\Omega}\subset K,
\]
which proves \eqref{eq:jump-support}.
Rectifiability follows from the rectifiability of $J_u$ and the
nullity of $K\setminus J_u$.

Finally, in the reduced case, restrict to a smooth $W\Subset\Omega$.
Every compactly supported comparison in $W$ extends by the
original pair outside $W$, with an unchanged collar.
Thus absolute minimality is preserved.
Reducedness is local and is preserved as well, and the
boundedness assertion follows from Step~2.
This also matches the locally Sobolev comparison class of the classical
theory. A competitor of finite energy has square-integrable gradient and,
on its bounded comparison region $V$, satisfies
\[
 \int_V |w|^2\leq
 2\int_V|w-g|^2+2|V|\,\|g\|_\infty^2.
\]
It has finite crack length there and agrees with the original pair on the
unchanged collar. Extension therefore gives an admissible competitor in
the original global $L^2$ and $W^{1,2}$ class. A competitor of infinite
energy cannot improve the original finite energy.
\end{proof}
